\chapter{Rational maps into group varieties}
\label{chap:Albanese}

This chapter develops the theory of rational maps needed for Milne's extension
theorem: \emph{a rational map from a nonsingular variety to an abelian variety
extends, uniquely, to a regular morphism}. The chapter proves the two halves that
do not mention the group structure of the target in an essential way --- the
calculus of rational maps (precomposition, pairing, function fields, the
difference map \(\Phi(x,y) = f(x)\,f(y)^{-1}\)), and the geometry of a smooth
variety (purity of polar loci, regular stalks, discrete valuation rings at
codimension-one points, Milne's codimension-\(\geq 2\) indeterminacy theorem) ---
and assembles them into the extension theorem, taking as an explicit hypothesis
the one remaining ingredient: Milne's Lemma~3.3, that the indeterminacy locus of
a rational map into a group variety is empty or of pure codimension one. That
lemma's statement is isolated as a named property
(\ref{def:milne33_indeterminacy}), and its principal substeps are proved here
(\ref{thm:difference_map_domain}, \ref{thm:difference_domain_reflection},
\ref{thm:difference_slice_witness}, \ref{thm:pole_purity},
\ref{thm:zero_locus_pure_codim_one}).

Throughout, a \emph{partial map} \((U, \varphi)\) from a scheme \(X\) to a scheme
\(Y\) is a dense open subset \(U \subseteq X\) together with a morphism
\(\varphi : U \to Y\); two partial maps are \emph{equivalent} if they agree on a
dense open subset contained in both domains. A \emph{rational map}
\(f : X \dashrightarrow Y\) is an equivalence class of partial maps. A point
\(x \in X\) is a \emph{point of definition} of \(f\) if some representative of
\(f\) is defined at \(x\); the points of definition form the dense open
\emph{domain of definition} \(\operatorname{Dom}(f)\). For a base scheme \(S\)
with \(X\), \(Y\) over \(S\), the rational map \(f\) is \emph{over \(S\)} when
some (equivalently, by shrinking, any small enough) representative is a morphism
over \(S\); when the structure morphisms are written \(s_X\), \(s_Y\), this is
the equation \(s_Y \circ f = s_X\) of rational maps \(X \dashrightarrow S\).
Composition of a rational map with a (total) morphism \(g : Y \to Z\) on the
left of the arrow, written \(g \circ f\), is represented by \((U, \varphi
\circ g)\) for any representative \((U, \varphi)\) of \(f\).

\section{Rational maps: precomposition, pairing, function fields}
\label{sec:ratmaps}

\begin{lemma}[Restriction square]
  \label{lem:ratmap_restrict_square}
  \lean{AlgebraicGeometry.Scheme.PartialMap.homOfLE_comp_morphismRestrict}
  \leanok
  Let \(p : W \to X\) be a morphism of schemes and \(V \subseteq U\) open subsets
  of \(X\). Write \(p_U : p^{-1}U \to U\) and \(p_V : p^{-1}V \to V\) for the
  restrictions of \(p\). Then the square of restrictions commutes:
  \[
    p_U \circ \iota_{p^{-1}V \subseteq p^{-1}U}
    \;=\; \iota_{V \subseteq U} \circ p_V
    \qquad\text{as morphisms } p^{-1}V \to U .
  \]
\end{lemma}
\begin{proof}
  \leanok
  The open immersion \(U \hookrightarrow X\) is a monomorphism, so it suffices to
  compare the two composites \(p^{-1}V \to U \hookrightarrow X\). Composing the
  left-hand side with \(U \hookrightarrow X\) gives the composite
  \(p^{-1}V \hookrightarrow p^{-1}U \hookrightarrow W \xrightarrow{p} X\), i.e.\
  the restriction of \(p\) to \(p^{-1}V\) followed by the inclusion into \(X\);
  composing the right-hand side with \(U \hookrightarrow X\) gives
  \(p^{-1}V \xrightarrow{p_V} V \hookrightarrow X\), which is the same morphism.
  Cancel the monomorphism.
\end{proof}

\begin{definition}[Precomposition of a rational map with an open morphism]
  \label{def:ratmap_precomp}
  \lean{AlgebraicGeometry.Scheme.PartialMap.precomp,
    AlgebraicGeometry.Scheme.PartialMap.precomp_equiv,
    AlgebraicGeometry.Scheme.RationalMap.precomp,
    AlgebraicGeometry.Scheme.RationalMap.precomp_toRationalMap}
  \uses{lem:ratmap_restrict_square}
  \leanok
  Let \(f : X \dashrightarrow Y\) be a rational map and \(p : W \to X\) a
  morphism whose underlying continuous map is \emph{open}. The composite
  \[
    f \circ p : W \dashrightarrow Y
  \]
  is the rational map represented by \((p^{-1}U,\ \varphi \circ p_U)\) for any
  representative \((U, \varphi)\) of \(f\), where \(p_U : p^{-1}U \to U\) is the
  restriction of \(p\).
\end{definition}
\begin{proof}
  \leanok
  Two things need proof: that \((p^{-1}U, \varphi \circ p_U)\) is a partial map,
  and that its class does not depend on the representative.

  \emph{Density.} If \(D \subseteq X\) is dense and \(p\) is an open map, then
  \(p^{-1}D\) is dense in \(W\): for any nonempty open \(O \subseteq W\) the
  image \(p(O)\) is a nonempty open subset of \(X\), hence meets \(D\), so there
  is \(w \in O\) with \(p(w) \in D\), i.e.\ \(O\) meets \(p^{-1}D\). Applied to
  \(D = U\), the preimage \(p^{-1}U\) is a dense open of \(W\).

  \emph{Independence of the representative.} Let \((U_1, \varphi_1)\) and
  \((U_2, \varphi_2)\) be equivalent representatives of \(f\), agreeing on a
  dense open \(V \subseteq U_1 \cap U_2\). Then \(p^{-1}V\) is dense (as above)
  and contained in \(p^{-1}U_1 \cap p^{-1}U_2\), and on \(p^{-1}V\) the two
  pulled-back morphisms agree: by the restriction square
  \ref{lem:ratmap_restrict_square}, restricting \(\varphi_i \circ p_{U_i}\) to
  \(p^{-1}V\) gives \(\varphi_i|_V \circ p_V\), and \(\varphi_1|_V =
  \varphi_2|_V\). Hence the two pullbacks are equivalent partial maps.
\end{proof}

\begin{lemma}[Precomposition preserves being over the base]
  \label{lem:ratmap_precomp_over}
  \lean{AlgebraicGeometry.Scheme.PartialMap.precomp_isOver}
  \uses{def:ratmap_precomp}
  \leanok
  Let \(S\) be a scheme, let \(W\), \(X\), \(Y\) be schemes over \(S\), let
  \((U, \varphi)\) be a partial map from \(X\) to \(Y\) whose morphism
  \(\varphi\) is a morphism over \(S\), and let \(p : W \to X\) be a morphism
  over \(S\) with open underlying map. Then the pulled-back partial map
  \((p^{-1}U, \varphi \circ p_U)\) is again over \(S\); in particular \(f \circ p\)
  is a rational map over \(S\) whenever \(f\) is.
\end{lemma}
\begin{proof}
  \leanok
  The restriction \(p_U : p^{-1}U \to U\) is a morphism over \(S\): composing it
  with the inclusion \(U \hookrightarrow X\) and the structure morphism of \(X\)
  gives the structure morphism of \(p^{-1}U\), because \(p\) is over \(S\) and
  open subschemes carry the restricted structure morphisms. A composite of
  morphisms over \(S\) is over \(S\).
\end{proof}

\begin{lemma}[Precomposition commutes with post-composition]
  \label{lem:ratmap_precomp_compHom}
  \lean{AlgebraicGeometry.Scheme.RationalMap.precomp_compHom}
  \uses{def:ratmap_precomp}
  \leanok
  For a rational map \(f : X \dashrightarrow Y\), a morphism \(p : W \to X\) with
  open underlying map, and a morphism \(g : Y \to Z\),
  \[
    g \circ (f \circ p) \;=\; (g \circ f) \circ p
    \qquad\text{as rational maps } W \dashrightarrow Z .
  \]
\end{lemma}
\begin{proof}
  \leanok
  Choose a representative \((U, \varphi)\) of \(f\). Both sides are represented
  by the partial map \(\bigl(p^{-1}U,\ g \circ (\varphi \circ p_U)\bigr)
  = \bigl(p^{-1}U,\ (g \circ \varphi) \circ p_U\bigr)\), by associativity of
  composition of morphisms.
\end{proof}

\begin{lemma}[Precomposing a morphism]
  \label{lem:ratmap_precomp_morphism}
  \lean{AlgebraicGeometry.Scheme.RationalMap.precomp_hom_toRationalMap}
  \uses{def:ratmap_precomp}
  \leanok
  Let \(g : X \to Y\) be a morphism, regarded as a rational map, and let
  \(p : W \to X\) have open underlying map. Then the precomposition of the
  rational map \(g\) with \(p\) is the rational map of the composite morphism
  \(g \circ p : W \to Y\).
\end{lemma}
\begin{proof}
  \leanok
  The morphism \(g\) is represented by the total partial map \((X, g)\); its
  pullback along \(p\) is \((p^{-1}X, g \circ p) = (W, g \circ p)\), the total
  partial map of the composite.
\end{proof}

\begin{lemma}[Precomposition enlarges the domain by the preimage]
  \label{lem:ratmap_le_domain_precomp}
  \lean{AlgebraicGeometry.Scheme.RationalMap.le_domain_precomp}
  \uses{def:ratmap_precomp}
  \leanok
  In the situation of \ref{def:ratmap_precomp},
  \[
    p^{-1}\bigl(\operatorname{Dom} f\bigr) \subseteq \operatorname{Dom}(f \circ p).
  \]
\end{lemma}
\begin{proof}
  \leanok
  Let \(w \in W\) with \(p(w) \in \operatorname{Dom} f\). By definition of the
  domain there is a representative \((U, \varphi)\) of \(f\) with \(p(w) \in U\).
  Its pullback \((p^{-1}U, \varphi \circ p_U)\) represents \(f \circ p\) and is
  defined at \(w\).
\end{proof}

\begin{lemma}[Post-composition only enlarges the domain]
  \label{lem:ratmap_le_domain_compHom}
  \lean{AlgebraicGeometry.Scheme.RationalMap.le_domain_compHom}
  \leanok
  For a rational map \(F : X \dashrightarrow Y\) and a morphism \(g : Y \to Z\),
  \[
    \operatorname{Dom}(F) \subseteq \operatorname{Dom}(g \circ F).
  \]
\end{lemma}
\begin{proof}
  \leanok
  A representative \((U, \varphi)\) of \(F\) defined at \(x\) yields the
  representative \((U, g \circ \varphi)\) of \(g \circ F\), with the same
  domain, defined at \(x\).
\end{proof}

\begin{lemma}[Associativity of post-composition]
  \label{lem:ratmap_compHom_compHom}
  \lean{AlgebraicGeometry.Scheme.RationalMap.compHom_compHom}
  \leanok
  For a rational map \(f : X \dashrightarrow Y\) and morphisms \(g : Y \to Z\),
  \(h : Z \to T\),
  \[
    h \circ (g \circ f) \;=\; (h \circ g) \circ f .
  \]
\end{lemma}
\begin{proof}
  \leanok
  Both sides are represented by \((U, h \circ g \circ \varphi)\) for any
  representative \((U, \varphi)\) of \(f\).
\end{proof}

The pairing of rational maps is built through the function field. Let \(X\) be
an \emph{integral} scheme with generic point \(\eta\) and function field
\(K(X) = \struct{X,\eta}\). Every open \(U \ni \eta\) receives a canonical
morphism \(\Spec K(X) \to U\) (the morphism of the local ring at \(\eta\)), so
each representative \((U, \varphi)\) of a rational map
\(f : X \dashrightarrow Y\) restricts to a morphism
\(\Spec K(X) \to Y\); two representatives agree on a dense open, which contains
\(\eta\), so this morphism is independent of the representative. We call it the
\emph{generic-point morphism} \(f_\eta : \Spec K(X) \to Y\) of \(f\).

\begin{theorem}[The function-field correspondence]
  \label{thm:ratmap_functionField_correspondence}
  \lean{AlgebraicGeometry.Scheme.RationalMap.fromFunctionField,
    AlgebraicGeometry.Scheme.RationalMap.ofFunctionField,
    AlgebraicGeometry.Scheme.RationalMap.equivFunctionField,
    AlgebraicGeometry.Scheme.RationalMap.eq_of_fromFunctionField_eq}
  \mathlibok
  \leanok
  Let \(X\) be an integral scheme. Two rational maps \(X \dashrightarrow Y\) are
  equal if and only if their generic-point morphisms agree. Moreover, if \(X\)
  and \(Y\) lie over a base \(S\) and the structure morphism \(s_Y\) is locally
  of finite type, every \(S\)-morphism \(\Spec K(X) \to Y\) is the generic-point
  morphism of a unique rational map \(X \dashrightarrow Y\) over \(S\): the
  finitely many coordinates of the image point, in an affine chart of \(Y\), are
  germs at \(\eta\) of functions defined on a common open of \(X\), so the
  morphism spreads out to a partial map.
\end{theorem}

\begin{lemma}[Generic-point morphism of a post-composition]
  \label{lem:ratmap_fromFunctionField_compHom}
  \lean{AlgebraicGeometry.Scheme.RationalMap.fromFunctionField_compHom}
  \uses{thm:ratmap_functionField_correspondence}
  \leanok
  Let \(X\) be irreducible, \(f : X \dashrightarrow Y\) a rational map and
  \(g : Y \to Z\) a morphism. Then \((g \circ f)_\eta = g \circ f_\eta\).
\end{lemma}
\begin{proof}
  \leanok
  Choose a representative \((U, \varphi)\) of \(f\); then \((U, g \circ \varphi)\)
  represents \(g \circ f\), and restricting to the generic point gives
  \((g \circ \varphi)_\eta = g \circ \varphi_\eta\) by associativity.
\end{proof}

\begin{lemma}[Generic-point morphism of a map over the base]
  \label{lem:ratmap_fromFunctionField_over}
  \lean{AlgebraicGeometry.Scheme.RationalMap.fromFunctionField_comp_structure}
  \uses{lem:ratmap_fromFunctionField_compHom}
  \leanok
  Let \(X\) be integral, let \(s_X : X \to S\), \(s_Y : Y \to S\) be structure
  morphisms and \(a : X \dashrightarrow Y\) a rational map over \(S\)
  (\(s_Y \circ a = s_X\) as rational maps). Then
  \[
    s_Y \circ a_\eta \;=\; s_X \circ c_X,
  \]
  where \(c_X : \Spec K(X) \to X\) is the canonical morphism.
\end{lemma}
\begin{proof}
  \leanok
  Apply \ref{lem:ratmap_fromFunctionField_compHom} with \(g = s_Y\):
  \((s_Y \circ a)_\eta = s_Y \circ a_\eta\). By hypothesis \(s_Y \circ a\) is the
  rational map of the morphism \(s_X\), whose generic-point morphism is the
  restriction of \(s_X\) to the generic point, i.e.\ \(s_X \circ c_X\).
\end{proof}

\begin{definition}[Pairing of rational maps into a fibre product]
  \label{def:ratmap_prod}
  \lean{AlgebraicGeometry.Scheme.RationalMap.prod,
    AlgebraicGeometry.Scheme.RationalMap.prod_fromFunctionField}
  \uses{thm:ratmap_functionField_correspondence, lem:ratmap_fromFunctionField_over}
  \leanok
  Let \(X\) be an integral scheme with structure morphism \(s_X : X \to S\), and
  let \(s_Y : Y \to S\), \(s_Z : Z \to S\) be locally of finite type. Let
  \(a : X \dashrightarrow Y\) and \(b : X \dashrightarrow Z\) be rational maps
  over \(S\). The \emph{pairing}
  \[
    (a, b) : X \dashrightarrow Y \times_S Z
  \]
  is the unique rational map over \(S\) whose generic-point morphism is the
  morphism \(\bigl(a_\eta, b_\eta\bigr) : \Spec K(X) \to Y \times_S Z\) induced
  by the universal property of the fibre product.
\end{definition}
\begin{proof}
  \leanok
  The two morphisms \(a_\eta\), \(b_\eta\) out of \(\Spec K(X)\) agree after
  composing with \(s_Y\), resp.\ \(s_Z\): both composites equal
  \(s_X \circ c_X\) by \ref{lem:ratmap_fromFunctionField_over}. Hence they
  factor through the fibre product. The structure morphism
  \(Y \times_S Z \to S\) is the composite of the projection to \(Y\) (a base
  change of \(s_Z\), hence locally of finite type) with \(s_Y\), so it is
  locally of finite type, and the function-field correspondence
  \ref{thm:ratmap_functionField_correspondence} produces the unique rational
  map over \(S\) with the prescribed generic-point morphism.
\end{proof}

\begin{lemma}[The projections of the pairing recover the factors]
  \label{lem:ratmap_prod_proj}
  \lean{AlgebraicGeometry.Scheme.RationalMap.prod_compHom_fst,
    AlgebraicGeometry.Scheme.RationalMap.prod_compHom_snd}
  \uses{def:ratmap_prod, lem:ratmap_fromFunctionField_compHom,
    thm:ratmap_functionField_correspondence}
  \leanok
  In the situation of \ref{def:ratmap_prod}, writing \(\mathrm{pr}_1\),
  \(\mathrm{pr}_2\) for the projections of \(Y \times_S Z\),
  \[
    \mathrm{pr}_1 \circ (a,b) = a
    \qquad\text{and}\qquad
    \mathrm{pr}_2 \circ (a,b) = b .
  \]
\end{lemma}
\begin{proof}
  \leanok
  Rational maps out of the integral \(X\) are determined by their generic-point
  morphisms (\ref{thm:ratmap_functionField_correspondence}). By
  \ref{lem:ratmap_fromFunctionField_compHom} the generic-point morphism of
  \(\mathrm{pr}_1 \circ (a,b)\) is \(\mathrm{pr}_1 \circ (a_\eta, b_\eta) =
  a_\eta\), and symmetrically for the second projection.
\end{proof}

\begin{lemma}[The pairing is over the base]
  \label{lem:ratmap_prod_over}
  \lean{AlgebraicGeometry.Scheme.RationalMap.prod_compHom_over}
  \uses{def:ratmap_prod, lem:ratmap_prod_proj, lem:ratmap_compHom_compHom}
  \leanok
  In the situation of \ref{def:ratmap_prod}, composing \((a,b)\) with the
  structure morphism \(s_Y \circ \mathrm{pr}_1\) of \(Y \times_S Z\) gives
  \(s_X\), as rational maps \(X \dashrightarrow S\).
\end{lemma}
\begin{proof}
  \leanok
  By associativity (\ref{lem:ratmap_compHom_compHom}) and
  \ref{lem:ratmap_prod_proj},
  \((s_Y \circ \mathrm{pr}_1) \circ (a,b) = s_Y \circ (\mathrm{pr}_1 \circ (a,b))
  = s_Y \circ a = s_X\).
\end{proof}

For a \emph{reduced} scheme \(X\) and a \emph{separated} scheme \(Y\), every
rational map \(f : X \dashrightarrow Y\) has a largest representative: the
representatives glue on overlaps (two morphisms into a separated scheme that
agree on a dense open of a reduced scheme agree, cf.\
\ref{lem:ratmap_hom_ext}), producing a representative defined on all of
\(\operatorname{Dom}(f)\).

\begin{definition}[The maximal representative]
  \label{def:ratmap_toPartialMap}
  \lean{AlgebraicGeometry.Scheme.RationalMap.toPartialMap,
    AlgebraicGeometry.Scheme.RationalMap.toRationalMap_toPartialMap}
  \mathlibok
  \leanok
  For \(X\) reduced and \(Y\) separated, the \emph{maximal representative} of a
  rational map \(f : X \dashrightarrow Y\) is the partial map with domain
  \(\operatorname{Dom}(f)\) obtained by gluing all representatives; its class
  is \(f\).
\end{definition}

\begin{definition}[The explicit pairing representative]
  \label{def:ratmap_pairPartialMap}
  \lean{AlgebraicGeometry.Scheme.RationalMap.pairPartialMap}
  \uses{def:ratmap_toPartialMap}
  \leanok
  Let \(X\) be reduced and let \(Y\), \(Z\), \(S\) be separated. Let
  \(a : X \dashrightarrow Y\), \(b : X \dashrightarrow Z\) be rational maps over
  \(S\), with maximal representatives \(\alpha\) on \(\operatorname{Dom}(a)\)
  and \(\beta\) on \(\operatorname{Dom}(b)\). The \emph{pairing representative}
  is the partial map from \(X\) to \(Y \times_S Z\) with domain
  \(\operatorname{Dom}(a) \cap \operatorname{Dom}(b)\) and morphism
  \((\alpha, \beta)\), the fibre-product pairing of the two restrictions.
\end{definition}
\begin{proof}
  \leanok
  The intersection of the dense open \(\operatorname{Dom}(a)\) with the dense
  open \(\operatorname{Dom}(b)\) is dense: a nonempty open \(O\) meets
  \(\operatorname{Dom}(a)\) in a nonempty open, which meets
  \(\operatorname{Dom}(b)\). The maximal representative of a rational map over
  the separated base \(S\) is a morphism over \(S\) (its restriction to a small
  enough dense open is, and the locus where the two composites to \(S\) agree is
  closed since \(S\) is separated, contains a dense open, and \(X\) is reduced);
  hence \(\alpha\) and \(\beta\) agree after composing to \(S\) and pair into
  the fibre product.
\end{proof}

\begin{lemma}[The pairing representative represents the pairing]
  \label{lem:ratmap_pairPartialMap_represents}
  \lean{AlgebraicGeometry.Scheme.RationalMap.pairPartialMap_toRationalMap}
  \uses{def:ratmap_pairPartialMap, def:ratmap_prod,
    thm:ratmap_functionField_correspondence}
  \leanok
  Let \(X\) be integral and reduced, \(Y\), \(Z\), \(S\) separated, \(s_Y\),
  \(s_Z\) locally of finite type, and \(a\), \(b\) as in
  \ref{def:ratmap_pairPartialMap}. Then the class of the pairing representative
  equals \((a, b)\).
\end{lemma}
\begin{proof}
  \leanok
  Both rational maps have the same generic-point morphism: restricting the
  pairing representative to the generic point and composing with each projection
  of \(Y \times_S Z\) gives the restriction of \(\alpha\), resp.\ \(\beta\), to
  the generic point, i.e.\ \(a_\eta\), resp.\ \(b_\eta\); and morphisms into a
  fibre product agree once their two projections do. Conclude by
  \ref{thm:ratmap_functionField_correspondence}.
\end{proof}

\begin{lemma}[The pairing is defined wherever both factors are]
  \label{lem:ratmap_le_domain_prod}
  \lean{AlgebraicGeometry.Scheme.RationalMap.le_domain_prod}
  \uses{lem:ratmap_pairPartialMap_represents}
  \leanok
  In the situation of \ref{lem:ratmap_pairPartialMap_represents},
  \[
    \operatorname{Dom}(a) \cap \operatorname{Dom}(b)
    \;\subseteq\; \operatorname{Dom}\bigl((a,b)\bigr).
  \]
\end{lemma}
\begin{proof}
  \leanok
  The pairing representative is a representative of \((a, b)\)
  (\ref{lem:ratmap_pairPartialMap_represents}) with domain
  \(\operatorname{Dom}(a) \cap \operatorname{Dom}(b)\), and the domain of a
  rational map contains the domain of each of its representatives.
\end{proof}

\begin{definition}[Germ pullback along a rational map]
  \label{def:ratmap_stalkPullback}
  \lean{AlgebraicGeometry.Scheme.RationalMap.stalkPullback}
  \uses{thm:ratmap_functionField_correspondence}
  \leanok
  Let \(X\) be an integral scheme and \(f : X \dashrightarrow Y\) a rational
  map. Let \(y \in Y\) be the image under \(f_\eta\) of the unique point of
  \(\Spec K(X)\). The \emph{germ pullback} of \(f\) is the local ring
  homomorphism
  \[
    f^\sharp : \struct{Y, y} \longrightarrow K(X)
  \]
  induced by \(f_\eta\) on the local ring at the image point.
\end{definition}

\begin{theorem}[A dominant rational map hits the generic point]
  \label{thm:ratmap_dominant_generic}
  \lean{AlgebraicGeometry.Scheme.RationalMap.fromFunctionField_base_eq_genericPoint}
  \uses{thm:ratmap_functionField_correspondence}
  \leanok
  Let \(X\) be integral, \(Y\) irreducible, and \(f : X \dashrightarrow Y\) a
  \emph{dominant} rational map (some, equivalently every, representative has
  dense image). Then the generic-point morphism \(f_\eta\) sends the unique
  point of \(\Spec K(X)\) to the generic point \(\eta_Y\) of \(Y\).
\end{theorem}
\begin{proof}
  \leanok
  Choose a representative \((U, \varphi)\) of \(f\); it is dominant. The open
  \(U\) is irreducible with generic point \(\eta = \eta_X\) (a dense open of an
  irreducible space is irreducible, and its generic point is that of the ambient
  space, which lies in every nonempty open). The canonical morphism
  \(\Spec K(X) \to U\) sends the unique point to \(\eta\), so
  \(f_\eta\) sends it to \(\varphi(\eta)\).

  It remains to see \(\varphi(\eta) = \eta_Y\). Since \(\{\eta\}\) is dense in
  \(U\) and \(\varphi\) is continuous,
  \(\varphi(U) = \varphi\bigl(\overline{\{\eta\}} \cap U\bigr) \subseteq
  \overline{\{\varphi(\eta)\}}\); taking closures and using dominance,
  \(\overline{\{\varphi(\eta)\}} \supseteq \overline{\varphi(U)} = Y\). So
  \(\varphi(\eta)\) is a generic point of the irreducible sober space \(Y\),
  hence equals \(\eta_Y\) by uniqueness of generic points.
\end{proof}

\begin{definition}[Function-field pullback of a dominant rational map]
  \label{def:ratmap_functionFieldPullback}
  \lean{AlgebraicGeometry.Scheme.RationalMap.functionFieldPullback}
  \uses{def:ratmap_stalkPullback, thm:ratmap_dominant_generic}
  \leanok
  Let \(X\) be integral, \(Y\) irreducible and \(f : X \dashrightarrow Y\)
  dominant. By \ref{thm:ratmap_dominant_generic} the source of the germ pullback
  \ref{def:ratmap_stalkPullback} is the local ring of \(Y\) at its generic
  point, i.e.\ the function field \(K(Y)\); the resulting field homomorphism
  \[
    f^* : K(Y) \longrightarrow K(X)
  \]
  is the \emph{function-field pullback} of \(f\).
\end{definition}

\section{The difference map}
\label{sec:difference_map}

Let \(\mathcal C\) be a cartesian monoidal category (finite products, with
\(\otimes\) the product and \(1\) the terminal object) and \(G\) a group object
of \(\mathcal C\), with multiplication \(\mu : G \otimes G \to G\), unit
\(e : 1 \to G\) and inversion \(\iota : G \to G\). For every object \(T\) the
set \(\Hom(T, G)\) is a group under \(u \cdot v = \mu \circ \langle u,
v\rangle\), with unit the constant \(e \circ t_T\) and inverse
\(u^{-1} = \iota \circ u\); precomposition with any morphism \(T' \to T\) is a
group homomorphism. We use this \emph{hom-group calculus} freely.

\begin{definition}[The difference morphism of a group object]
  \label{def:grp_diff}
  \lean{CategoryTheory.GrpObj.diff}
  \leanok
  For a group object \(G\) in a cartesian monoidal category, the
  \emph{difference morphism} is
  \[
    \delta \;=\; \mathrm{pr}_1 \cdot \mathrm{pr}_2^{-1}
    \;:\; G \otimes G \longrightarrow G,
  \]
  the quotient, in the hom-group \(\Hom(G \otimes G, G)\), of the two
  projections; on \(T\)-points it is \((g, h) \mapsto g h^{-1}\).
\end{definition}

\begin{lemma}[Reconstruction identity]
  \label{lem:grp_diff_reconstruct}
  \lean{CategoryTheory.GrpObj.lift_diff_lift_mul}
  \uses{def:grp_diff}
  \leanok
  For \(T\)-points \(a, b : T \to G\) of a group object \(G\),
  \[
    \mu \circ \bigl\langle\, \delta \circ \langle a, b\rangle,\ b \,\bigr\rangle
    \;=\; a .
  \]
  In pointwise notation: \((a b^{-1})\, b = a\).
\end{lemma}
\begin{proof}
  \leanok
  Precomposition with \(\langle a, b\rangle\) is compatible with the hom-group
  operations, so
  \(\delta \circ \langle a, b\rangle
  = (\mathrm{pr}_1 \circ \langle a,b\rangle)\cdot(\mathrm{pr}_2 \circ
  \langle a,b\rangle)^{-1} = a \cdot b^{-1}\) in \(\Hom(T, G)\). Then
  \(\mu \circ \langle a b^{-1}, b\rangle = (a b^{-1}) \cdot b = a\) by the group
  axioms of the hom-group.
\end{proof}

Now fix a field \(k\). We work in the category of schemes over \(k\); the
categorical product of \(X\) and \(Y\) over \(k\) is the fibre product
\(X \times_k Y\). A \emph{group scheme} over \(k\) is a group object \(G\) in
this category.

\begin{definition}[The difference and multiplication morphisms of a group scheme]
  \label{def:grp_scheme_diff}
  \lean{AlgebraicGeometry.grpObjDiffLeft,
    AlgebraicGeometry.grpObjMulLeft,
    AlgebraicGeometry.grpObjDiffLeft_comp_hom}
  \uses{def:grp_diff}
  \leanok
  For a group scheme \(G\) over \(k\), write
  \[
    \delta_G : G \times_k G \longrightarrow G,
    \qquad
    m_G : G \times_k G \longrightarrow G
  \]
  for the underlying scheme morphisms of the difference morphism
  \ref{def:grp_diff} and of the multiplication. Both are morphisms over \(k\):
  writing \(s_G\) for the structure morphism of \(G\), one has
  \(s_G \circ \delta_G = s_G \circ \mathrm{pr}_1\), the structure morphism of
  \(G \times_k G\).
\end{definition}

\begin{lemma}[Scheme-level reconstruction identity]
  \label{lem:grp_scheme_diff_reconstruct}
  \lean{AlgebraicGeometry.pullback_lift_diff_lift_mul}
  \uses{def:grp_scheme_diff, lem:grp_diff_reconstruct}
  \leanok
  Let \(G\) be a group scheme over \(k\), \(T\) a scheme with a morphism
  \(t : T \to \Spec k\), and \(A, B : T \to G\) two morphisms with
  \(s_G \circ A = s_G \circ B = t\), where \(s_G\) is the structure morphism.
  Then
  \[
    m_G \circ \bigl\langle\, \delta_G \circ \langle A, B\rangle,\ B \,\bigr\rangle
    \;=\; A ,
  \]
  where the pairings are taken into the fibre product \(G \times_k G\).
\end{lemma}
\begin{proof}
  \leanok
  The morphisms \(A\), \(B\) are morphisms over \(k\) from the object
  \((T, t)\), i.e.\ \(T\)-points of the group object \(G\) in the category of
  schemes over \(k\); the fibre-product pairings are the categorical pairings
  there. The identity is \ref{lem:grp_diff_reconstruct} for these points, read
  off on underlying scheme morphisms (the forgetful functor from schemes over
  \(k\) to schemes preserves the relevant composites and sends categorical
  pairings to fibre-product pairings).
\end{proof}

\begin{lemma}[The projections of a smooth self-product are open]
  \label{lem:self_product_open}
  \lean{AlgebraicGeometry.isOpenMap_pullback_fst_self,
    AlgebraicGeometry.isOpenMap_pullback_snd_self}
  \leanok
  Let \(X\) be a scheme smooth over \(k\). Then both projections
  \(\mathrm{pr}_1, \mathrm{pr}_2 : X \times_k X \to X\) are open maps.
\end{lemma}
\begin{proof}
  \leanok
  Each projection is a base change of the smooth structure morphism
  \(X \to \Spec k\), hence smooth; smooth morphisms are flat and locally of
  finite presentation, and such morphisms are open.
\end{proof}

\begin{definition}[Milne's difference rational map]
  \label{def:difference_map}
  \lean{AlgebraicGeometry.Scheme.RationalMap.differenceRationalMap}
  \uses{def:ratmap_precomp, lem:self_product_open, def:ratmap_prod,
    def:grp_scheme_diff, lem:ratmap_precomp_over, lem:ratmap_precomp_compHom,
    lem:ratmap_precomp_morphism}
  \leanok
  \source{abelian-varieties:page-0023}
  Let \(X\) be a scheme smooth over \(k\) whose self-product \(X \times_k X\) is
  integral, and let \(G\) be a group scheme locally of finite type over \(k\).
  For a rational map \(f : X \dashrightarrow G\) over \(k\), the
  \emph{difference rational map}
  \[
    \Phi \;:\; X \times_k X \dashrightarrow G,
    \qquad
    \Phi(x, y) = f(x)\, f(y)^{-1},
  \]
  is the composite \(\delta_G \circ \bigl(f \circ \mathrm{pr}_1,\ f \circ
  \mathrm{pr}_2\bigr)\): precompose \(f\) with the two (open, by
  \ref{lem:self_product_open}) projections, pair the results into
  \(G \times_k G\) (\ref{def:ratmap_prod}; both factors are over \(k\) by
  \ref{lem:ratmap_precomp_over}, and the pairing needs the integrality of
  \(X \times_k X\)), and post-compose with the difference morphism
  \(\delta_G\).
\end{definition}

\begin{lemma}[The difference map is over the base]
  \label{lem:difference_map_over}
  \lean{AlgebraicGeometry.Scheme.RationalMap.differenceRationalMap_compHom_over}
  \uses{def:difference_map, lem:ratmap_compHom_compHom, lem:ratmap_prod_over,
    def:grp_scheme_diff}
  \leanok
  In the situation of \ref{def:difference_map}, \(\Phi\) is a rational map over
  \(k\): composing \(\Phi\) with the structure morphism of \(G\) gives the
  structure morphism of \(X \times_k X\), as rational maps.
\end{lemma}
\begin{proof}
  \leanok
  By associativity of post-composition (\ref{lem:ratmap_compHom_compHom}) and
  since \(\delta_G\) is a morphism over \(k\) (\ref{def:grp_scheme_diff}),
  composing \(\Phi\) with the structure morphism of \(G\) equals composing the
  pairing with the structure morphism of \(G \times_k G\), which is the
  structure morphism of \(X \times_k X\) by \ref{lem:ratmap_prod_over}.
\end{proof}

\begin{definition}[The explicit pairing representative of the difference map]
  \label{def:difference_pairing}
  \lean{AlgebraicGeometry.Scheme.RationalMap.precompDiffPairing}
  \uses{def:ratmap_toPartialMap, def:ratmap_precomp, lem:self_product_open,
    lem:ratmap_precomp_over}
  \leanok
  Let \(X\) be smooth over \(k\) with \(X\) reduced, and \(G\) a separated group
  scheme over \(k\). Let \(f : X \dashrightarrow G\) be a rational map over
  \(k\), with maximal representative \(\varphi\) on
  \(U = \operatorname{Dom}(f)\). The \emph{pairing representative} is the
  partial map from \(X \times_k X\) to \(G \times_k G\) with domain
  \(\mathrm{pr}_1^{-1}(U) \cap \mathrm{pr}_2^{-1}(U)\) and morphism the
  fibre-product pairing
  \(\bigl(\varphi \circ \mathrm{pr}_1,\ \varphi \circ \mathrm{pr}_2\bigr)\)
  of the two pulled-back restrictions.
\end{definition}
\begin{proof}
  \leanok
  The two preimages are dense opens (preimages of the dense open \(U\) under
  the open projections, as in \ref{def:ratmap_precomp}), and the intersection of
  two dense opens is dense. Both pulled-back morphisms are morphisms over
  \(k\) (\ref{lem:ratmap_precomp_over}, using that both projections are
  morphisms over \(k\)), so they pair into the fibre product
  \(G \times_k G\).
\end{proof}

\begin{lemma}[The pairing representative represents the pairing of the pullbacks]
  \label{lem:difference_pairing_represents}
  \lean{AlgebraicGeometry.Scheme.RationalMap.precompDiffPairing_toRationalMap}
  \uses{def:difference_pairing, def:ratmap_prod,
    thm:ratmap_functionField_correspondence}
  \leanok
  In the situation of \ref{def:difference_pairing}, assume moreover that
  \(X \times_k X\) is integral and \(G\) is locally of finite type over \(k\).
  Then the class of the pairing representative equals the pairing
  \(\bigl(f \circ \mathrm{pr}_1,\ f \circ \mathrm{pr}_2\bigr)\) of
  \ref{def:difference_map}.
\end{lemma}
\begin{proof}
  \leanok
  As in \ref{lem:ratmap_pairPartialMap_represents}: both sides are rational maps
  out of the integral scheme \(X \times_k X\), and their generic-point morphisms
  agree component-wise --- composing with each projection of \(G \times_k G\)
  gives the generic-point morphism of \(f \circ \mathrm{pr}_i\), computed either
  from the pulled-back maximal representative or from the pairing. Conclude by
  \ref{thm:ratmap_functionField_correspondence}.
\end{proof}

\begin{theorem}[The difference map is defined on
    \(\operatorname{Dom}(f) \times \operatorname{Dom}(f)\)]
  \label{thm:difference_map_domain}
  \lean{AlgebraicGeometry.Scheme.RationalMap.le_domain_differenceRationalMap}
  \uses{def:difference_map, lem:difference_pairing_represents,
    lem:ratmap_le_domain_compHom}
  \leanok
  \source{abelian-varieties:page-0023}
  In the situation of \ref{def:difference_map}, with \(X\) reduced and \(G\)
  separated,
  \[
    \mathrm{pr}_1^{-1}\bigl(\operatorname{Dom} f\bigr) \cap
    \mathrm{pr}_2^{-1}\bigl(\operatorname{Dom} f\bigr)
    \;\subseteq\;
    \operatorname{Dom}(\Phi).
  \]
  In particular \(\Phi\) is defined at every point \((x, y)\) with
  \(x, y \in \operatorname{Dom}(f)\).
\end{theorem}
\begin{proof}
  \leanok
  Post-composition with \(\delta_G\) only enlarges the domain
  (\ref{lem:ratmap_le_domain_compHom}), so it suffices to bound the domain of
  the pairing; and the pairing representative
  (\ref{lem:difference_pairing_represents}) is a representative of the pairing
  defined exactly on
  \(\mathrm{pr}_1^{-1}(\operatorname{Dom} f) \cap
  \mathrm{pr}_2^{-1}(\operatorname{Dom} f)\).
\end{proof}

\begin{theorem}[Reconstruction of \(f \circ \mathrm{pr}_1\) from the difference map]
  \label{thm:difference_reconstruct}
  \lean{AlgebraicGeometry.Scheme.RationalMap.reconstruct_precomp_fst}
  \uses{def:difference_map, lem:difference_map_over, def:ratmap_prod,
    lem:grp_scheme_diff_reconstruct, lem:ratmap_fromFunctionField_compHom,
    thm:ratmap_functionField_correspondence, lem:ratmap_fromFunctionField_over}
  \leanok
  \source{abelian-varieties:page-0023}
  In the situation of \ref{def:difference_map},
  \[
    m_G \circ \bigl(\Phi,\ f \circ \mathrm{pr}_2\bigr)
    \;=\; f \circ \mathrm{pr}_1
    \qquad\text{as rational maps } X \times_k X \dashrightarrow G,
  \]
  where the pairing is formed as in \ref{def:ratmap_prod} (both entries are over
  \(k\), by \ref{lem:difference_map_over}). This is Milne's formula
  \(f(x) = \Phi(x, u) \cdot f(u)\).
\end{theorem}
\begin{proof}
  \leanok
  Both sides are rational maps out of the integral scheme \(X \times_k X\); we
  compare generic-point morphisms
  (\ref{thm:ratmap_functionField_correspondence}). Write
  \(A = (f \circ \mathrm{pr}_1)_\eta\) and \(B = (f \circ \mathrm{pr}_2)_\eta\),
  two morphisms \(\Spec K(X \times_k X) \to G\) with equal composites to
  \(\Spec k\) (\ref{lem:ratmap_fromFunctionField_over}). Unwinding
  \ref{def:difference_map} and \ref{def:ratmap_prod} with
  \ref{lem:ratmap_fromFunctionField_compHom}, the generic-point morphism of the
  left-hand side is
  \(m_G \circ \bigl\langle \delta_G \circ \langle A, B\rangle,\ B\bigr\rangle\),
  which equals \(A\) by the scheme-level reconstruction identity
  \ref{lem:grp_scheme_diff_reconstruct}.
\end{proof}

\begin{theorem}[Domain reflection through the difference map]
  \label{thm:difference_domain_reflection}
  \lean{AlgebraicGeometry.Scheme.RationalMap.le_domain_precomp_fst_of_difference}
  \uses{thm:difference_reconstruct, lem:ratmap_le_domain_compHom,
    lem:ratmap_le_domain_prod, lem:ratmap_le_domain_precomp}
  \leanok
  \source{abelian-varieties:page-0023}
  In the situation of \ref{thm:difference_reconstruct}, with \(X\) reduced and
  \(G\) separated,
  \[
    \operatorname{Dom}(\Phi) \,\cap\,
    \mathrm{pr}_2^{-1}\bigl(\operatorname{Dom} f\bigr)
    \;\subseteq\;
    \operatorname{Dom}\bigl(f \circ \mathrm{pr}_1\bigr).
  \]
  At a point where \(\Phi\) is defined and whose second coordinate is a point of
  definition of \(f\), the reconstruction \(f(x) = \Phi(x, u)\, f(u)\) exhibits
  \(f \circ \mathrm{pr}_1\) as defined.
\end{theorem}
\begin{proof}
  \leanok
  Rewrite \(f \circ \mathrm{pr}_1 = m_G \circ (\Phi, f \circ \mathrm{pr}_2)\)
  by \ref{thm:difference_reconstruct}. The right-hand side is defined wherever
  the pairing is (\ref{lem:ratmap_le_domain_compHom}), the pairing is defined on
  \(\operatorname{Dom}(\Phi) \cap \operatorname{Dom}(f \circ \mathrm{pr}_2)\)
  (\ref{lem:ratmap_le_domain_prod}), and
  \(\mathrm{pr}_2^{-1}(\operatorname{Dom} f) \subseteq
  \operatorname{Dom}(f \circ \mathrm{pr}_2)\)
  (\ref{lem:ratmap_le_domain_precomp}).
\end{proof}

\section{Purity of the polar locus}
\label{sec:pole_purity}

This section proves that on an integral, locally Noetherian scheme whose local
rings are all regular, the locus where a fixed element of the function field
fails to be regular is of \emph{pure codimension one}: every point of the locus
is a specialisation of a codimension-one point of the locus
(\ref{thm:pole_purity}). The classical route runs through normality and the
Krull-domain intersection \(R = \bigcap_{\mathrm{ht}\,\mathfrak p = 1}
R_{\mathfrak p}\); the proof given here is elementary, replacing the
integral-closedness input by a two-element ``swap'' argument valid in any
regular local ring of dimension \(\geq 2\).

\begin{lemma}[Swap lemma]
  \label{lem:swap_lemma}
  \lean{mem_span_singleton_of_swap_pair}
  \leanok
  Let \(A\) be a domain, and let \(u, v, b, x \in A\) with \(u \neq 0\) and
  \(v\) regular modulo \((u)\) (i.e.\ \(v s \in (u)\) implies \(s \in (u)\)).
  If \(u x \in (b)\) and \(v x \in (b)\), then \(x \in (b)\).
\end{lemma}
\begin{proof}
  \leanok
  If \(b = 0\), then \(u x = 0\) and \(u \neq 0\) force \(x = 0 \in (0)\).
  Assume \(b \neq 0\) and write \(u x = b s\), \(v x = b t\). Then
  \[
    b (u t) \;=\; u (v x) \;=\; v (u x) \;=\; b (v s),
  \]
  and cancelling \(b\) (a nonzero element of a domain) gives \(u t = v s\).
  Thus \(v s \in (u)\), so \(s \in (u)\) by regularity of \(v\) modulo
  \((u)\): write \(s = u s'\). Then \(u x = b u s'\), and cancelling \(u\)
  gives \(x = b s' \in (b)\).
\end{proof}

\begin{theorem}[Swap pairs exist in regular local rings of dimension \(\geq 2\)]
  \label{thm:swap_pair_exists}
  \lean{IsRegularLocalRing.exists_swap_pair_of_two_le_ringKrullDim}
  \uses{lem:notMemSq_witness, lem:regular_quotient, thm:regular_isDomain}
  \leanok
  Let \((A, \mathfrak m)\) be a regular local (Noetherian) ring of Krull
  dimension \(\geq 2\). Then there are \(u, v \in \mathfrak m\) with
  \(u \neq 0\) and \(v\) regular modulo \((u)\).
\end{theorem}
\begin{proof}
  \leanok
  In a regular local ring the minimal number of generators
  \(\mu(\mathfrak m)\) equals the Krull dimension, so
  \(\mu(\mathfrak m) \geq 2\). Choose
  \(u \in \mathfrak m \setminus \mathfrak m^2\)
  (\ref{lem:notMemSq_witness}); then \(u \neq 0\) since
  \(0 \in \mathfrak m^2\). The quotient \(A' = A/(u)\) is again a regular local
  ring, with \(\mu(\mathfrak m') = \mu(\mathfrak m) - 1 \geq 1\)
  (\ref{lem:regular_quotient}); being regular local it is a domain
  (\ref{thm:regular_isDomain}), and \(\mathfrak m' \neq 0\) because
  \(\mu(\mathfrak m') \geq 1\). Pick a nonzero \(\bar v \in \mathfrak m'\) and
  lift it to \(v \in A\); then \(v \in \mathfrak m\) (a unit would have a unit
  image, but \(\bar v \in \mathfrak m'\)), and \(v\) is regular modulo
  \((u)\): if \(v s \in (u)\) then \(\bar v \bar s = 0\) in the domain \(A'\)
  with \(\bar v \neq 0\), so \(\bar s = 0\), i.e.\ \(s \in (u)\).
\end{proof}

\begin{theorem}[Height bound for colon primes of principal ideals]
  \label{thm:colon_height_le_one}
  \lean{Ideal.height_le_one_of_colon_span_singleton}
  \uses{lem:swap_lemma, thm:swap_pair_exists}
  \leanok
  Let \(A\) be a Noetherian domain, \(b, y \in A\), and let
  \(Q = \bigl((b) : y\bigr)\) be a prime ideal. If the localisation \(A_Q\) is a
  regular local ring, then \(\mathrm{ht}(Q) \leq 1\).
\end{theorem}
\begin{proof}
  \leanok
  Suppose \(\mathrm{ht}(Q) \geq 2\); then \(\dim A_Q = \mathrm{ht}(Q) \geq 2\)
  (the primes of \(A_Q\) are the primes of \(A\) inside \(Q\)), and \(A_Q\) is a
  domain (a localisation of a domain). Write \(\varphi : A \to A_Q\) for the
  localisation map, injective since \(A\) is a domain.

  \emph{Step 1: \(\varphi(y) \notin (\varphi(b))\).} Otherwise
  \(\varphi(y) = \varphi(b) \cdot c/s\) with \(s \notin Q\); cross-multiplying
  and using injectivity gives \(s y \in (b)\) in \(A\), i.e.\
  \(s \in \bigl((b) : y\bigr) = Q\), a contradiction.

  \emph{Step 2: the maximal ideal of \(A_Q\) is the colon
  \(\bigl((\varphi b) : \varphi y\bigr)\).} For \(r \in Q\) we have
  \(r y \in (b)\), so \(\varphi(r) \varphi(y) \in (\varphi b)\); since the
  maximal ideal of \(A_Q\) is generated by \(\varphi(Q)\), this gives
  \(\mathfrak m_{A_Q} \subseteq \bigl((\varphi b) : \varphi y\bigr)\).
  Conversely the colon is a proper ideal by Step 1 (it does not contain \(1\)),
  hence is contained in the maximal ideal.

  \emph{Step 3: contradiction by the swap lemma.} By
  \ref{thm:swap_pair_exists} there is a swap pair
  \(u, v \in \mathfrak m_{A_Q}\), \(u \neq 0\), \(v\) regular modulo \((u)\).
  By Step 2, \(u \cdot \varphi(y) \in (\varphi b)\) and
  \(v \cdot \varphi(y) \in (\varphi b)\); the swap lemma \ref{lem:swap_lemma}
  gives \(\varphi(y) \in (\varphi b)\), contradicting Step 1.
\end{proof}

\begin{theorem}[Existence of a height-one pole prime]
  \label{thm:pole_prime_height_one}
  \lean{exists_height_one_prime_colon_le}
  \uses{thm:colon_height_le_one}
  \leanok
  Let \(R\) be a Noetherian domain all of whose localisations at primes are
  regular local rings. Let \(a, b \in R\) with \(b \neq 0\), and let
  \(\mathfrak p\) be a prime containing the denominator ideal
  \(\bigl((b) : a\bigr)\). Then there is a prime
  \(\mathfrak q \subseteq \mathfrak p\) of height exactly \(1\) with
  \(\bigl((b) : a\bigr) \subseteq \mathfrak q\).
\end{theorem}
\begin{proof}
  \leanok
  Work in the localisation \(A = R_{\mathfrak p}\), a Noetherian domain, with
  injective localisation map \(\varphi : R \to A\); primes of \(A\) correspond
  to primes of \(R\) contained in \(\mathfrak p\), compatibly with inclusions
  and heights.

  \emph{Step 1: \(\varphi(a) \notin (\varphi b)\).} Otherwise, cross-multiplying
  as in Step~1 of \ref{thm:colon_height_le_one}, some \(s \notin \mathfrak p\)
  has \(s a \in (b)\), i.e.\ \(s \in \bigl((b) : a\bigr) \subseteq \mathfrak p\)
  --- a contradiction. Hence the class of \(\varphi(a)\) in \(A/(\varphi b)\) is
  a nonzero element of a finitely generated module over the Noetherian ring
  \(A\).

  \emph{Step 2: an associated prime.} The annihilator of a nonzero element of a
  finitely generated module over a Noetherian ring is contained in an associated
  prime of that module, and an associated prime is by definition the annihilator
  of some element. So there is a prime \(Q \subseteq A\), the annihilator of the
  class of some \(y \in A\) in \(A/(\varphi b)\) --- that is,
  \(Q = \bigl((\varphi b) : y\bigr)\) --- with
  \(\bigl((\varphi b) : \varphi a\bigr) \subseteq Q\).

  \emph{Step 3: \(\mathrm{ht}(Q) = 1\).} The localisation \(A_Q\) is a
  localisation of \(R\) at the prime \(Q \cap R\), hence regular by hypothesis;
  \ref{thm:colon_height_le_one} gives \(\mathrm{ht}(Q) \leq 1\). Moreover
  \(\varphi(b) \in Q\) (as \(\varphi(b)\, y \in (\varphi b)\)) and
  \(\varphi(b) \neq 0\), so \(Q \neq 0\), and in a domain a nonzero prime has
  height \(\geq 1\) (the chain \(0 \subsetneq Q\)). Hence
  \(\mathrm{ht}(Q) = 1\).

  \emph{Step 4: descent to \(R\).} Let \(\mathfrak q = Q \cap R\), a prime of
  \(R\) contained in \(\mathfrak p\) (as \(Q\) is contained in the maximal ideal
  of \(A\), whose contraction is \(\mathfrak p\)). The correspondence of primes
  preserves heights, so \(\mathrm{ht}(\mathfrak q) = 1\). Finally, for
  \(r \in \bigl((b) : a\bigr)\): \(\varphi(r)\varphi(a) \in (\varphi b)\), so
  \(\varphi(r) \in \bigl((\varphi b) : \varphi a\bigr) \subseteq Q\), i.e.\
  \(r \in \mathfrak q\).
\end{proof}

\begin{theorem}[Purity of the polar locus]
  \label{thm:pole_purity}
  \lean{AlgebraicGeometry.Scheme.exists_specializes_coheight_eq_one_of_notMem_stalk_range}
  \uses{thm:pole_prime_height_one, thm:stalk_dim_coheight}
  \leanok
  \source{abelian-varieties:page-0024}
  Let \(X\) be an integral, locally Noetherian scheme all of whose stalks
  \(\struct{X,x}\) are regular local rings. Let \(h \in K(X)\) and let
  \(P \in X\) be a point at which \(h\) is \emph{not regular}, i.e.\ \(h\) does
  not lie in the image of \(\struct{X,P} \to K(X)\). Then there is a point
  \(z \in X\) with
  \[
    z \rightsquigarrow P, \qquad
    \operatorname{coheight}(z) = 1, \qquad
    h \notin \operatorname{im}\bigl(\struct{X,z} \to K(X)\bigr),
  \]
  where \(z \rightsquigarrow P\) means \(P \in \overline{\{z\}}\).
\end{theorem}
\begin{proof}
  \leanok
  Choose an affine open \(U \ni P\), \(U \cong \Spec R\) with
  \(R = \Gamma(U, \struct U)\); \(R\) is a Noetherian domain, its fraction field
  is \(K(X)\), and for every point \(w \in U\) with prime \(\mathfrak w\) the
  stalk \(\struct{X,w}\) is the localisation \(R_{\mathfrak w}\). In particular
  every localisation of \(R\) at a prime is a regular local ring, by the
  hypothesis on stalks. Let \(\mathfrak p\) be the prime of \(P\) and write
  \(h = a/b\) with \(a, b \in R\), \(b \neq 0\).

  \emph{Step 1 (regularity as a colon condition).} For a prime
  \(\mathfrak q \subseteq R\),
  \[
    \frac{a}{b} \in \operatorname{im}\bigl(R_{\mathfrak q} \to K(X)\bigr)
    \iff
    \bigl((b) : a\bigr) \not\subseteq \mathfrak q .
  \]
  (\(\Rightarrow\)) If \(a/b = c/s\) with \(s \notin \mathfrak q\), then
  \(s a = b c\) in \(R\) (cross-multiply in the fraction field and use
  injectivity of \(R \to K(X)\)), so \(s \in \bigl((b) : a\bigr)\) and the colon
  is not contained in \(\mathfrak q\). (\(\Leftarrow\)) If
  \(s \in \bigl((b) : a\bigr) \setminus \mathfrak q\), write \(s a = b c\); then
  \(a/b = c/s \in R_{\mathfrak q}\).

  \emph{Step 2 (the height-one prime).} Since \(h\) is not regular at \(P\),
  Step 1 (applied to \(\mathfrak q = \mathfrak p\), whose localisation is
  \(\struct{X,P}\)) gives \(\bigl((b) : a\bigr) \subseteq \mathfrak p\). By
  \ref{thm:pole_prime_height_one} there is a prime
  \(\mathfrak q \subseteq \mathfrak p\) of height \(1\) with
  \(\bigl((b) : a\bigr) \subseteq \mathfrak q\).

  \emph{Step 3 (the point).} Let \(z \in U\) be the point of \(\mathfrak q\).
  Then \(z \rightsquigarrow P\): inclusions of primes correspond to
  specialisations in \(\Spec R \cong U\), and specialisations in the open
  subspace \(U\) are specialisations in \(X\). The stalk \(\struct{X,z}\) is
  \(R_{\mathfrak q}\), of Krull dimension \(\mathrm{ht}(\mathfrak q) = 1\), so
  \(\operatorname{coheight}(z) = 1\) by \ref{thm:stalk_dim_coheight}. Finally
  \(h\) is not regular at \(z\) by Step 1, since
  \(\bigl((b) : a\bigr) \subseteq \mathfrak q\).
\end{proof}

\section{Two substeps of Milne's Lemma 3.3}
\label{sec:milne33_substeps}

Milne's Lemma 3.3 asserts that the indeterminacy locus of a rational map from a
nonsingular variety to a group variety is empty or of pure codimension one
(the statement is isolated as \ref{def:milne33_indeterminacy} below). This
section proves two of its substeps: the topological slice argument producing
the auxiliary point \(u\) of Milne's formula \(f(x) = \Phi(x,u)\,f(u)\), and
the Krull-theoretic fact that the zero locus of a regular function is of pure
codimension one.

\begin{theorem}[Slice witness on the self-product]
  \label{thm:fibre_slice_witness}
  \lean{AlgebraicGeometry.exists_snd_mem_of_fst_eq_of_mem}
  \uses{lem:self_product_open}
  \leanok
  \source{abelian-varieties:page-0023}
  Let \(X\) be a scheme, smooth and geometrically irreducible over the field
  \(k\). Let \(V \subseteq X \times_k X\) be open, \(D \subseteq X\) a dense
  open, \(x \in X\) a point, and suppose some point \(p_0\) of the fibre
  \(\mathrm{pr}_1^{-1}\{x\}\) lies in \(V\). Then there is a point \(p\) with
  \[
    \mathrm{pr}_1(p) = x, \qquad p \in V, \qquad \mathrm{pr}_2(p) \in D .
  \]
\end{theorem}
\begin{proof}
  \leanok
  \emph{The fibre is irreducible.} The projection \(\mathrm{pr}_1\) is a base
  change of the structure morphism of \(X\), hence geometrically irreducible
  (geometric irreducibility is stable under base change) and open
  (\ref{lem:self_product_open}); the preimage of an irreducible subset under an
  open morphism with irreducible geometric fibres is irreducible. Applied to
  \(\{x\}\), the fibre \(F = \mathrm{pr}_1^{-1}\{x\}\) is irreducible.

  \emph{Both opens meet the fibre.} \(V\) meets \(F\) at \(p_0\). For
  \(\mathrm{pr}_2^{-1}(D)\): pick \(u \in D\) (dense opens of the nonempty
  \(X\) are nonempty). Over the one-point base \(\Spec k\) there is a point
  \(q\) of \(X \times_k X\) with \(\mathrm{pr}_1(q) = x\) and
  \(\mathrm{pr}_2(q) = u\): the residue fields \(\kappa(x)\) and \(\kappa(u)\)
  are extensions of \(k\), their tensor product \(\kappa(x) \otimes_k
  \kappa(u)\) is a nonzero ring, and any prime of it gives such a point. Thus
  \(q \in F \cap \mathrm{pr}_2^{-1}(D)\).

  \emph{Conclusion.} An irreducible set is preirreducible: two open sets each
  meeting \(F\) meet \(F\) simultaneously. Applied to \(V\) and
  \(\mathrm{pr}_2^{-1}(D)\) this produces the desired \(p \in F \cap V \cap
  \mathrm{pr}_2^{-1}(D)\).
\end{proof}

\begin{theorem}[Slice witness for the difference map]
  \label{thm:difference_slice_witness}
  \lean{AlgebraicGeometry.Scheme.RationalMap.exists_mem_domain_precomp_fst_of_differenceRationalMap}
  \uses{thm:fibre_slice_witness, thm:difference_domain_reflection}
  \leanok
  \source{abelian-varieties:page-0023}
  In the situation of \ref{def:difference_map}, with \(X\) moreover
  geometrically irreducible over \(k\), \(X\) reduced and \(G\) separated:
  if the difference map \(\Phi\) is defined at some point \(p_0\) with
  \(\mathrm{pr}_1(p_0) = x\), then there is a point \(p\) with
  \(\mathrm{pr}_1(p) = x\) at which \(f \circ \mathrm{pr}_1\) is defined.
\end{theorem}
\begin{proof}
  \leanok
  Apply \ref{thm:fibre_slice_witness} with \(V = \operatorname{Dom}(\Phi)\)
  (open, met by the fibre at \(p_0\)) and \(D = \operatorname{Dom}(f)\)
  (dense open): this yields \(p\) over \(x\) with \(p \in
  \operatorname{Dom}(\Phi)\) and \(\mathrm{pr}_2(p) \in \operatorname{Dom}(f)\).
  By the domain reflection \ref{thm:difference_domain_reflection}, \(f \circ
  \mathrm{pr}_1\) is defined at \(p\).
\end{proof}

\begin{theorem}[Height-one transport of Krull's Hauptidealsatz]
  \label{thm:krull_height_one_transport}
  \lean{Ideal.exists_height_one_prime_mem_le}
  \leanok
  \source{stacks-project}
  Let \(R\) be a Noetherian domain, \(t \in R\) nonzero, and \(\mathfrak p\) a
  prime containing \(t\). Then there is a prime
  \(\mathfrak q \subseteq \mathfrak p\) of height exactly \(1\) with
  \(t \in \mathfrak q\).
\end{theorem}
\begin{proof}
  \leanok
  The ideal \((t)\) is contained in \(\mathfrak p\), so among the primes between
  \((t)\) and \(\mathfrak p\) there is one, say \(\mathfrak q\), minimal over
  \((t)\) (the set of primes containing \((t)\) and contained in
  \(\mathfrak p\) is nonempty and closed under intersections of chains). By
  Krull's principal ideal theorem, a prime minimal over a principal ideal in a
  Noetherian ring has height \(\leq 1\). And \(t \in \mathfrak q\) with
  \(t \neq 0\) gives \(\mathfrak q \neq 0\), so in the domain \(R\) the chain
  \(0 \subsetneq \mathfrak q\) gives \(\mathrm{ht}(\mathfrak q) \geq 1\).
\end{proof}

\begin{lemma}[Regularity propagates to generisations]
  \label{lem:regular_at_generisation}
  \lean{AlgebraicGeometry.range_algebraMap_stalk_le_of_specializes}
  \leanok
  Let \(X\) be an irreducible scheme and \(z \rightsquigarrow P\) a
  specialisation. Then
  \[
    \operatorname{im}\bigl(\struct{X,P} \to K(X)\bigr)
    \subseteq
    \operatorname{im}\bigl(\struct{X,z} \to K(X)\bigr):
  \]
  an element of the function field regular at \(P\) is regular at every
  generisation \(z\) of \(P\).
\end{lemma}
\begin{proof}
  \leanok
  Since \(z\) is a generisation of \(P\), every open set containing \(P\)
  contains \(z\); taking germs therefore defines a canonical homomorphism
  \(\struct{X,P} \to \struct{X,z}\), compatible with the two maps into the
  common stalk at the generic point. Hence the image of \(\struct{X,P}\) in
  \(K(X)\) factors through that of \(\struct{X,z}\).
\end{proof}

\begin{theorem}[The zero locus of a regular function is pure codimension one]
  \label{thm:zero_locus_pure_codim_one}
  \lean{AlgebraicGeometry.Scheme.exists_specializes_coheight_eq_one_of_mem_maximalIdeal}
  \uses{thm:pole_purity, lem:regular_at_generisation}
  \leanok
  \source{abelian-varieties:page-0024}
  Let \(X\) be an integral, locally Noetherian scheme with regular stalks. Let
  \(P \in X\) and let \(s \in \struct{X,P}\) lie in the maximal ideal, with
  nonzero image \(t \in K(X)\). Then there is a point \(z \rightsquigarrow P\)
  with \(\operatorname{coheight}(z) = 1\) and an element
  \(s' \in \struct{X,z}\), again in the maximal ideal, with the same image
  \(t\) in \(K(X)\).
\end{theorem}
\begin{proof}
  \leanok
  \emph{Step A: \(t^{-1}\) is not regular at \(P\).} Suppose
  \(t^{-1}\) were the image of some \(u \in \struct{X,P}\). Then \(u s \in
  \mathfrak m_P\), so \(1 - u s\) is a unit of the local ring
  \(\struct{X,P}\); but its image in \(K(X)\) is
  \(1 - t^{-1} t = 0\), and a ring homomorphism into the nonzero ring \(K(X)\)
  sends units to units --- a contradiction.

  \emph{Step B: purity.} By \ref{thm:pole_purity} applied to \(t^{-1}\), there
  is \(z \rightsquigarrow P\) with \(\operatorname{coheight}(z) = 1\) and
  \(t^{-1}\) not regular at \(z\).

  \emph{Step C: \(t\) is regular at \(z\).} By
  \ref{lem:regular_at_generisation}, \(t\), being regular at \(P\), is regular
  at the generisation \(z\): there is \(s' \in \struct{X,z}\) with image
  \(t\).

  \emph{Step D: \(s'\) lies in the maximal ideal.} Otherwise \(s'\) is a unit
  of \(\struct{X,z}\), and the image of \(s'^{-1}\) in \(K(X)\) is
  \(t^{-1}\), making \(t^{-1}\) regular at \(z\) --- contradicting Step B.
\end{proof}

\section{Smooth varieties: charts, reducedness, regular stalks,
  discrete valuation rings}
\label{sec:smooth_stalks}

This section records the local structure of a scheme smooth over a field:
standard-smooth charts, freeness and rank of the K\"ahler differentials through
localisation, the cotangent computation at rational points, reducedness,
regularity of all stalks, integrality of the self-product, and the discrete
valuation rings at codimension-one points.

\begin{theorem}[Standard-smooth charts]
  \label{thm:smooth_chart}
  \lean{AlgebraicGeometry.Scheme.exists_isStandardSmooth_at_of_smooth}
  \leanok
  \source{stacks-project}
  Let \(\bar k\) be an algebraically closed field and \(X\) a scheme smooth
  over \(\bar k\). Every point \(z \in X\) has an affine open neighbourhood
  \(V\), lying over an affine open \(U \subseteq \Spec \bar k\), such that the
  induced ring map \(\Gamma(U) \to \Gamma(V)\) is \emph{standard smooth}: a
  quotient \(\Gamma(U)[x_1, \dots, x_n]/(f_1, \dots, f_c)\) whose Jacobian
  determinant \(\det(\partial f_i/\partial x_j)_{i,j \leq c}\) is invertible.
\end{theorem}
\begin{proof}
  \leanok
  Smoothness is local on source and target, and a smooth ring map is, locally
  on the source, standard smooth (the local structure theorem for smooth
  morphisms, Stacks 00T7): choose an affine chart of \(X\) around \(z\) and
  shrink it accordingly.
\end{proof}

\begin{lemma}[K\"ahler differentials of a standard-smooth algebra are free]
  \label{lem:ss_kaehler_free}
  \lean{AlgebraicGeometry.Scheme.module_free_kaehlerDifferential_of_isStandardSmooth}
  \leanok
  \source{stacks-project}
  Let \(S\) be a standard-smooth algebra over a ring \(R\). Then
  \(\Omega_{S/R}\) is a free \(S\)-module.
\end{lemma}
\begin{proof}
  \leanok
  Write \(S = R[x_1, \dots, x_n]/(f_1, \dots, f_c)\) with invertible Jacobian
  minor. In the conormal sequence of this presentation the classes of
  \(f_1, \dots, f_c\) map to a basis of the summand
  \(\bigoplus_{i \leq c} S\,dx_i\) (invertibility of the Jacobian determinant),
  so \(\Omega_{S/R}\) is free with basis the images of
  \(dx_{c+1}, \dots, dx_n\) (Stacks 00T7).
\end{proof}

\begin{lemma}[Freeness of K\"ahler differentials localises]
  \label{lem:kaehler_free_localization}
  \lean{AlgebraicGeometry.Scheme.module_free_kaehlerDifferential_localization}
  \leanok
  \source{stacks-project}
  Let \(S\) be an \(R\)-algebra with \(\Omega_{S/R}\) a free \(S\)-module, and
  let \(S_M\) be a localisation of \(S\) at a multiplicative set \(M\). Then
  \(\Omega_{S_M/R}\) is a free \(S_M\)-module.
\end{lemma}
\begin{proof}
  \leanok
  K\"ahler differentials localise:
  \(\Omega_{S_M/R} \cong S_M \otimes_S \Omega_{S/R}\) (Stacks 00RT). The base
  change of a free module is free.
\end{proof}

\begin{lemma}[Rank of the localised K\"ahler differentials]
  \label{lem:kaehler_rank_localization}
  \lean{AlgebraicGeometry.Scheme.rank_kaehlerDifferential_localization_eq_relativeDimension}
  \uses{lem:ss_kaehler_free, lem:kaehler_free_localization}
  \leanok
  \source{stacks-project}
  Let \(S\) be a nonzero standard-smooth \(R\)-algebra of relative dimension
  \(n\), and \(S_M\) a nonzero localisation of \(S\). Then \(\Omega_{S_M/R}\)
  is free of rank \(n\) over \(S_M\).
\end{lemma}
\begin{proof}
  \leanok
  \(\Omega_{S/R}\) is free of rank \(n\) over \(S\): the basis exhibited in
  \ref{lem:ss_kaehler_free} has \(n - c\) elements for a presentation on \(n'\)
  variables and \(c\) relations with \(n = n' - c\) --- the relative dimension.
  Localisation is a base change, so \(\Omega_{S_M/R} \cong S_M \otimes_S
  \Omega_{S/R}\) (\ref{lem:kaehler_free_localization}) is free of the same rank
  (well-defined since \(S_M \neq 0\)).
\end{proof}

\begin{theorem}[Cotangent dimension at a rational point]
  \label{thm:cotangent_finrank_rational}
  \lean{AlgebraicGeometry.Scheme.finrank_cotangentSpace_of_bijective_algebraMap_residue}
  \leanok
  \source{stacks-project}
  Let \(R\) be a ring and \((S, \mathfrak m, \kappa)\) a local \(R\)-algebra,
  formally smooth over \(R\), with \(\Omega_{S/R}\) free of rank
  \(n \in \mathbb{N}\). If the composite \(R \to S \to \kappa\) is
  \emph{bijective}, then
  \[
    \dim_\kappa \mathfrak m/\mathfrak m^2 = n .
  \]
\end{theorem}
\begin{proof}
  \leanok
  Since \(R \to \kappa\) is an isomorphism of rings, \(\kappa\) is formally
  smooth over \(R\), and \(\Omega_{\kappa/R} = 0\) (every \(R\)-derivation of
  \(\kappa\) kills the image of \(R\), which is everything).

  Consider the conormal map of the surjection \(S \twoheadrightarrow \kappa\)
  with kernel \(\mathfrak m\) (Stacks 00RU):
  \[
    c : \mathfrak m/\mathfrak m^2 \longrightarrow
    \kappa \otimes_S \Omega_{S/R},
    \qquad c(\bar f) = df \otimes 1 .
  \]
  \emph{Injectivity.} Because \(\kappa\) is formally smooth over \(R\) and
  \(S \to \kappa\) is surjective, the conormal map admits a retraction (formal
  smoothness of the quotient provides compatible liftings to the square-zero
  extensions \(S/\mathfrak m^2 \twoheadrightarrow \kappa\), splitting \(c\));
  in particular \(c\) is injective.

  \emph{Surjectivity.} The conormal sequence
  \(\mathfrak m/\mathfrak m^2 \to \kappa \otimes_S \Omega_{S/R} \to
  \Omega_{\kappa/R} \to 0\) is exact, and \(\Omega_{\kappa/R} = 0\), so \(c\)
  is surjective.

  Hence \(c\) is an isomorphism of \(\kappa\)-vector spaces, and the right-hand
  side, the base change of a free rank-\(n\) module, has dimension \(n\).
\end{proof}

\begin{lemma}[Sections of the point over a field]
  \label{lem:gamma_spec_field}
  \lean{AlgebraicGeometry.Scheme.gammaSpecField_ringEquiv}
  \leanok
  Let \(k\) be a field and \(U\) a nonempty open of \(\Spec k\). Then
  \(\Gamma(\Spec k, U) \cong k\) as rings.
\end{lemma}
\begin{proof}
  \leanok
  \(\Spec k\) has a single point, so the only nonempty open is the whole space,
  whose ring of sections is \(k\).
\end{proof}

\begin{theorem}[Regularity at a rational closed point of a standard-smooth
    algebra]
  \label{thm:ss_regular_rational}
  \lean{AlgebraicGeometry.Scheme.isRegularLocalRing_localization_of_isStandardSmooth_of_bijective_residue}
  \uses{lem:kaehler_free_localization, lem:kaehler_rank_localization,
    thm:cotangent_finrank_rational, lem:standard_smooth_dim_lower_localized,
    lem:regular_of_cotangent_le_dim}
  \leanok
  Let \(k\) be a field, \(S\) a nonzero standard-smooth \(k\)-algebra,
  \(\mathfrak m \subseteq S\) a maximal ideal, and \(S_{\mathfrak m}\) a
  localisation of \(S\) at \(\mathfrak m\). If the residue map
  \(k \to \kappa(\mathfrak m)\) is bijective, then \(S_{\mathfrak m}\) is a
  regular local ring.
\end{theorem}
\begin{proof}
  \leanok
  \(S\) is of finite type over \(k\), hence Noetherian, and so is the
  localisation \(S_{\mathfrak m}\). Fix \(n\) with \(S\) standard smooth of
  relative dimension \(n\) (every standard-smooth algebra has a presentation
  with a well-defined relative dimension). By
  \ref{lem:kaehler_free_localization} and \ref{lem:kaehler_rank_localization},
  \(\Omega_{S_{\mathfrak m}/k}\) is free of rank \(n\). The localisation of a
  (formally) smooth algebra is formally smooth, so
  \ref{thm:cotangent_finrank_rational} applies and gives
  \(\dim_\kappa \mathfrak m/\mathfrak m^2 = n\) for the maximal ideal of
  \(S_{\mathfrak m}\). On the other side,
  \(\dim S_{\mathfrak m} \geq n\) by the height lower bound at maximal ideals
  (\ref{lem:standard_smooth_dim_lower_localized}). So
  \(\dim_\kappa \mathfrak m/\mathfrak m^2 = n \leq \dim S_{\mathfrak m}\),
  and \ref{lem:regular_of_cotangent_le_dim} concludes.
\end{proof}

\begin{theorem}[Regularity at every closed point over an algebraically closed
    field]
  \label{thm:ss_regular_closed}
  \lean{AlgebraicGeometry.Scheme.isRegularLocalRing_localizationAtPrime_of_isStandardSmooth_of_isAlgClosed}
  \uses{thm:ss_regular_rational}
  \leanok
  Let \(\bar k\) be an algebraically closed field and \(S\) a nonzero
  standard-smooth \(\bar k\)-algebra. For every maximal ideal
  \(\mathfrak m \subseteq S\), the localisation \(S_{\mathfrak m}\) is a
  regular local ring.
\end{theorem}
\begin{proof}
  \leanok
  By Zariski's lemma the residue field \(S/\mathfrak m\), a finitely generated
  \(\bar k\)-algebra that is a field, is a finite --- hence algebraic ---
  extension of \(\bar k\); as \(\bar k\) is algebraically closed, the residue
  map \(\bar k \to \kappa(\mathfrak m)\) is bijective. Apply
  \ref{thm:ss_regular_rational}.
\end{proof}

\begin{theorem}[Standard-smooth algebras over an algebraically closed field are
    reduced]
  \label{thm:ss_reduced}
  \lean{AlgebraicGeometry.Scheme.isReduced_of_isStandardSmooth_of_isAlgClosed}
  \uses{thm:ss_regular_closed, thm:regular_isDomain}
  \leanok
  \source{stacks-project}
  Let \(\bar k\) be an algebraically closed field and \(S\) a standard-smooth
  \(\bar k\)-algebra. Then \(S\) is reduced.
\end{theorem}
\begin{proof}
  \leanok
  The zero ring is reduced; assume \(S \neq 0\). A ring is reduced as soon as
  its localisations at all maximal ideals are reduced (a nonzero nilpotent
  remains nonzero in some such localisation, since its annihilator is a proper
  ideal). Each localisation \(S_{\mathfrak m}\) is regular local
  (\ref{thm:ss_regular_closed}), hence a domain
  (\ref{thm:regular_isDomain}), hence reduced.
\end{proof}

\begin{theorem}[A scheme smooth over an algebraically closed field is reduced]
  \label{thm:smooth_scheme_reduced}
  \lean{AlgebraicGeometry.Scheme.isReduced_of_smooth_of_isAlgClosed}
  \uses{thm:smooth_chart, lem:gamma_spec_field, thm:ss_reduced}
  \leanok
  \source{stacks-project}
  Let \(\bar k\) be an algebraically closed field and \(X\) a scheme smooth
  over \(\bar k\). Then \(X\) is reduced.
\end{theorem}
\begin{proof}
  \leanok
  Reducedness is checked on stalks. Fix \(z \in X\) and choose a
  standard-smooth chart \(V \ni z\) over \(U \subseteq \Spec \bar k\)
  (\ref{thm:smooth_chart}); identifying \(\Gamma(U) \cong \bar k\)
  (\ref{lem:gamma_spec_field}, \(U\) is nonempty as it contains the image of
  \(z\)) makes \(\Gamma(V)\) a standard-smooth \(\bar k\)-algebra, which is
  reduced by \ref{thm:ss_reduced}. The stalk \(\struct{X,z}\) is a localisation
  of \(\Gamma(V)\), and localisations of reduced rings are reduced.
\end{proof}

\begin{theorem}[Regularity of all stalks of a smooth variety]
  \label{thm:smooth_stalk_regular}
  \lean{AlgebraicGeometry.Scheme.isRegularLocalRing_stalk_of_smooth}
  \uses{thm:smooth_chart, lem:gamma_spec_field, thm:smooth_prime_regularity}
  \leanok
  \source{stacks-project}
  Let \(\bar k\) be an algebraically closed field and \(X\) an integral scheme,
  smooth, geometrically irreducible, separated and locally of finite type over
  \(\bar k\). Then the stalk \(\struct{X,z}\) at \emph{every} point \(z \in X\)
  (closed or not) is a regular local ring.
\end{theorem}
\begin{proof}
  \leanok
  Choose a standard-smooth chart \(V \ni z\) over \(U\)
  (\ref{thm:smooth_chart}) and identify \(\Gamma(U) \cong \bar k\)
  (\ref{lem:gamma_spec_field}), making \(\Gamma(V)\) a nonzero standard-smooth
  \(\bar k\)-algebra. The stalk \(\struct{X,z}\) is the localisation of
  \(\Gamma(V)\) at the prime of \(z\). Since \(\bar k\) is perfect, the
  smooth-prime regularity theorem \ref{thm:smooth_prime_regularity} applies at
  the (arbitrary) prime of \(z\) and gives regularity of the localisation.
\end{proof}

\begin{theorem}[Integrality of the self-product]
  \label{thm:self_product_integral}
  \lean{AlgebraicGeometry.Scheme.isIntegral_pullback_self}
  \uses{thm:smooth_scheme_reduced}
  \leanok
  Let \(\bar k\) be an algebraically closed field and \(X\) an integral scheme,
  smooth and geometrically irreducible over \(\bar k\). Then the self fibre
  product \(X \times_{\bar k} X\) is an integral scheme.
\end{theorem}
\begin{proof}
  \leanok
  The structure morphism of \(X \times_{\bar k} X\) is the composite of the
  projection \(\mathrm{pr}_1\) (a base change of the smooth structure morphism
  of \(X\), hence smooth) with the structure morphism of \(X\); so
  \(X \times_{\bar k} X\) is smooth over \(\bar k\) and therefore reduced
  (\ref{thm:smooth_scheme_reduced}). It is also irreducible: the structure
  morphism of \(X\) is geometrically irreducible and universally open (being
  smooth), and the fibre product of the irreducible \(X\) with a scheme whose
  structure morphism is geometrically irreducible and universally open is
  irreducible (the projection is open with irreducible fibres onto an
  irreducible base). Reduced and irreducible is integral.
\end{proof}

\begin{theorem}[Matsumura's regular-sequence criterion]
  \label{thm:matsumura_regular_sequence}
  \lean{matsumura_isRegular_of_linearIndependent_cotangent}
  \uses{thm:regular_isDomain, lem:regular_quotient}
  \leanok
  Let \((A, \mathfrak m, \kappa)\) be a regular local Noetherian ring and
  \(f_1, \dots, f_c \in \mathfrak m\) elements whose classes in
  \(\mathfrak m/\mathfrak m^2\) are \(\kappa\)-linearly independent. Then
  \(f_1, \dots, f_c\) is a regular sequence in \(A\).
\end{theorem}
\begin{proof}
  \leanok
  Induction on \(c\); the empty sequence is regular. Let \(c \geq 1\).

  \emph{The head is a nonzerodivisor.} \(A\) is a domain
  (\ref{thm:regular_isDomain}), and \([f_1] \neq 0\) in
  \(\mathfrak m/\mathfrak m^2\) (linear independence), so
  \(f_1 \notin \mathfrak m^2\); in particular \(f_1 \neq 0\), hence \(f_1\) is
  a nonzerodivisor.

  \emph{The quotient is regular local.} Since
  \(f_1 \in \mathfrak m \setminus \mathfrak m^2\), the quotient
  \(A' = A/(f_1)\) is again a regular local ring, with minimal generator count
  \(\mu(\mathfrak m') = \mu(\mathfrak m) - 1\) (\ref{lem:regular_quotient};
  note \(\mu(\mathfrak m) \geq 1\) because \(f_1 \in \mathfrak m \setminus
  \mathfrak m^2\) forces \(\mathfrak m \neq 0\)).

  \emph{Descent of linear independence.} Let \(\pi : \mathfrak m/\mathfrak m^2
  \to \mathfrak m'/\mathfrak m'^2\) be the map induced by \(A \twoheadrightarrow
  A'\) (well defined: \(\mathfrak m\) maps onto \(\mathfrak m'\) and
  \(\mathfrak m^2\) into \(\mathfrak m'^2\); the residue fields of \(A\) and
  \(A'\) agree). \(\pi\) is surjective, and its kernel is the \(\kappa\)-line
  spanned by \([f_1]\): an element \(x \in \mathfrak m\) has
  \(\pi([x]) = 0\) exactly when \(x \in (f_1) + \mathfrak m^2\), i.e.\
  \(x = a f_1 + m\) with \(m \in \mathfrak m^2\); writing \(a = a_0 + a'\) with
  \(a' \in \mathfrak m\) gives \([x] = \bar a_0 [f_1]\), and conversely.
  Now suppose \(\sum_{i \geq 2} \bar g_i\, \pi([f_i]) = 0\) with
  \(\bar g_i \in \kappa\). Lifting the coefficients to \(A\) and applying the
  kernel computation, \(\sum_{i \geq 2} \bar g_i [f_i] \in \kappa [f_1]\), so
  some relation \(\bar b [f_1] - \sum_{i \geq 2} \bar g_i [f_i] = 0\) holds in
  \(\mathfrak m / \mathfrak m^2\); linear independence of
  \([f_1], \dots, [f_c]\) forces all coefficients, in particular every
  \(\bar g_i\), to vanish. Hence the classes of \(f_2, \dots, f_c\) in
  \(\mathfrak m'/\mathfrak m'^2\) are \(\kappa\)-linearly independent.

  \emph{Conclusion.} By the inductive hypothesis applied to the regular local
  Noetherian ring \(A'\), the images of \(f_2, \dots, f_c\) form a regular
  sequence in \(A'\). Together with the nonzerodivisor \(f_1\) this makes
  \(f_1, f_2, \dots, f_c\) a regular sequence in \(A\): the successive
  quotients \(A/(f_1, \dots, f_i) \cong A'/(\bar f_2, \dots, \bar f_i)\) agree,
  so each \(f_{i+1}\) acts as a nonzerodivisor on the corresponding quotient,
  and the total quotient \(A/(f_1, \dots, f_c) \neq 0\) since \(A'\) is local
  and the images lie in its maximal ideal.
\end{proof}

\begin{theorem}[Smooth and codimension one implies a discrete valuation ring]
  \label{thm:smooth_codim_one_dvr}
  \lean{AlgebraicGeometry.Scheme.localRing_dvr_of_codim_one}
  \uses{thm:smooth_stalk_regular, thm:stalk_dim_coheight, thm:regular_isDomain}
  \leanok
  \source{stacks-project}
  Let \(\bar k\) be an algebraically closed field and \(X\) an integral scheme,
  smooth, geometrically irreducible, separated and locally of finite type over
  \(\bar k\). For every point \(z \in X\) with
  \(\operatorname{coheight}(z) = 1\), the stalk \(\struct{X,z}\) is a discrete
  valuation ring.
\end{theorem}
\begin{proof}
  \leanok
  The stalk \(A = \struct{X,z}\) is Noetherian (\(X\) is locally of finite
  type over the field \(\bar k\), hence locally Noetherian) and regular local
  (\ref{thm:smooth_stalk_regular}), of Krull dimension
  \(\dim A = \operatorname{coheight}(z) = 1\)
  (\ref{thm:stalk_dim_coheight}). Regularity in dimension one says
  \(\dim_\kappa \mathfrak m/\mathfrak m^2 = 1\), so by Nakayama the maximal
  ideal is principal, generated by any lift of a basis vector, and
  \(\mathfrak m \neq 0\) (else \(\dim A = 0\)). A regular local ring is a
  domain (\ref{thm:regular_isDomain}); and a Noetherian local domain whose
  maximal ideal is principal and nonzero is a discrete valuation ring
  (Stacks 00PD).
\end{proof}

\section{The indeterminacy locus and Milne's Theorem 3.1}
\label{sec:indeterminacy}

\begin{definition}[Indeterminacy locus]
  \label{def:indeterminacy_locus}
  \lean{AlgebraicGeometry.Scheme.RationalMap.indeterminacyLocus}
  \leanok
  \source{abelian-varieties:page-0022}
  The \emph{indeterminacy locus} of a rational map
  \(f : X \dashrightarrow Y\) is the closed subset
  \[
    Z(f) \;=\; X \setminus \operatorname{Dom}(f).
  \]
  The map \(f\) is represented by a (unique, for \(X\) reduced and \(Y\)
  separated) total morphism if and only if \(Z(f) = \emptyset\).
\end{definition}

\begin{definition}[Codimension-one-indeterminacy-free rational maps]
  \label{def:codim_one_free}
  \lean{AlgebraicGeometry.Scheme.RationalMap.CodimOneFree}
  \uses{def:indeterminacy_locus}
  \leanok
  A rational map \(f : X \dashrightarrow Y\) is
  \emph{codimension-one-indeterminacy-free} if every point \(x \in X\) with
  \(\operatorname{coheight}(x) = 1\) lies in \(\operatorname{Dom}(f)\) ---
  equivalently, if no codimension-one point of \(X\) lies in \(Z(f)\).
\end{definition}

\begin{definition}[The Milne 3.3 indeterminacy property]
  \label{def:milne33_indeterminacy}
  \lean{AlgebraicGeometry.Scheme.RationalMap.Milne33Indeterminacy}
  \uses{def:indeterminacy_locus}
  \leanok
  \source{abelian-varieties:page-0023}
  A rational map \(f : X \dashrightarrow G\) \emph{has the Milne 3.3
  indeterminacy property} if its indeterminacy locus is empty or of pure
  codimension one, in the following precise sense: either
  \(Z(f) = \emptyset\), or every \(x \in Z(f)\) satisfies
  \[
    \exists\, z \in Z(f), \quad \operatorname{coheight}(z) = 1
    \ \text{ and }\ x \in \overline{\{z\}} .
  \]
  (The requirement \(z \in Z(f)\) is essential: on a variety every non-generic
  point is a specialisation of \emph{some} codimension-one point, so without it
  the second alternative would be nearly vacuous, and in particular too weak to
  combine with the codimension-\(\geq 2\) bound of
  \ref{thm:indeterminacy_codimGe2}.)

  Milne's Lemma 3.3 asserts that every rational map over an algebraically
  closed field \(\bar k\) from a nonsingular \(\bar k\)-variety to a group
  \(\bar k\)-variety has this property. That assertion is \emph{not} proved in
  this chapter; the results proved here towards it are
  \ref{thm:difference_map_domain}, \ref{thm:difference_reconstruct},
  \ref{thm:difference_domain_reflection}, \ref{thm:difference_slice_witness},
  \ref{thm:pole_purity} and \ref{thm:zero_locus_pure_codim_one}. Every
  consumer below takes the property as an explicit hypothesis.
\end{definition}

\begin{theorem}[Milne's Theorem 3.1: indeterminacy in codimension \(\geq 2\)]
  \label{thm:indeterminacy_codimGe2}
  \lean{AlgebraicGeometry.Scheme.RationalMap.indeterminacy_codimGe2_of_smooth_of_complete}
  \uses{def:indeterminacy_locus, thm:smooth_codim_one_dvr,
    thm:ratmap_functionField_correspondence,
    lem:ratmap_fromFunctionField_compHom}
  \leanok
  \source{abelian-varieties:page-0022, abelian-varieties:page-0023}
  Let \(\bar k\) be an algebraically closed field, \(X\) a nonsingular
  \(\bar k\)-variety (an integral scheme, smooth, geometrically irreducible,
  separated and locally of finite type over \(\bar k\)) and \(Y\) a complete
  \(\bar k\)-variety (an integral scheme, proper and geometrically irreducible
  over \(\bar k\)). Let \(f : X \dashrightarrow Y\) be a rational map over
  \(\bar k\). Then every point of \(Z(f)\) has coheight at least \(2\):
  \[
    \forall\, z \in Z(f), \quad 2 \leq \operatorname{coheight}(z).
  \]
\end{theorem}
\begin{proof}
  \leanok
  Let \(z \in Z(f)\) and suppose \(\operatorname{coheight}(z) \leq 1\).

  \emph{Case \(\operatorname{coheight}(z) = 0\).} Then \(z\) admits no strict
  generisation, i.e.\ it is maximal in the specialisation preorder. The generic
  point \(\eta\) of the irreducible sober space \(X\) generizes every point, so
  maximality forces \(z = \eta\). But the generic point lies in every nonempty
  open subset, in particular in the dense open \(\operatorname{Dom}(f)\) ---
  contradicting \(z \in Z(f)\).

  \emph{Case \(\operatorname{coheight}(z) = 1\).} The stalk
  \(\struct{X,z}\) is a discrete valuation ring
  (\ref{thm:smooth_codim_one_dvr}) whose fraction field is the function field
  \(K(X)\) (the stalk of an integral scheme injects into the stalk at the
  generic point, and at a codimension-one point of a Noetherian integral scheme
  the fraction field of the stalk is \(K(X)\)). Consider the commutative
  square
  \[
    \begin{array}{ccc}
      \Spec K(X) & \xrightarrow{\ f_\eta\ } & Y \\
      \downarrow & & \downarrow \\
      \Spec \struct{X,z} & \longrightarrow & \Spec \bar k
    \end{array}
  \]
  whose left leg is induced by the inclusion \(\struct{X,z} \subseteq K(X)\)
  and whose bottom is the structure morphism of \(\Spec \struct{X,z}\) (the
  square commutes because \(f\) is a rational map over \(\bar k\): composing
  \(f_\eta\) with the structure morphism of \(Y\) gives the structure morphism
  of \(\Spec K(X)\), by \ref{lem:ratmap_fromFunctionField_compHom} applied to
  the over-\(\bar k\) equation). Since \(Y \to \Spec \bar k\) is proper and
  \(\struct{X,z}\) is a valuation ring with fraction field \(K(X)\), the
  valuative criterion of properness produces a lift
  \(L : \Spec \struct{X,z} \to Y\) with
  \(L|_{\Spec K(X)} = f_\eta\), compatible with the structure morphisms.

  Because \(Y\) is locally of finite type over \(\bar k\) (properness includes
  finite type), the morphism \(L\), defined at the level of the local ring at
  \(z\), spreads out to a morphism \(g : V \to Y\) on an open neighbourhood
  \(V \ni z\): in an affine chart of \(Y\) around the image of the closed point
  of \(\Spec \struct{X,z}\), the finitely many coordinate germs of \(L\) are
  defined on a common neighbourhood of \(z\). The partial map \((V, g)\) --- its
  domain is dense because \(X\) is irreducible --- has generic-point morphism
  equal to the restriction of \(L\) to the generic point, i.e.\ \(f_\eta\); by
  the function-field correspondence
  (\ref{thm:ratmap_functionField_correspondence}) its class is \(f\). Hence
  \(z \in \operatorname{Dom}(f)\) --- contradicting \(z \in Z(f)\).
\end{proof}

\begin{corollary}[Rational maps to complete varieties are codimension-one-free]
  \label{cor:codim_one_free_complete}
  \lean{AlgebraicGeometry.Scheme.RationalMap.codimOneFree_of_smooth_of_complete}
  \uses{thm:indeterminacy_codimGe2, def:codim_one_free}
  \leanok
  \source{abelian-varieties:page-0022}
  In the situation of \ref{thm:indeterminacy_codimGe2}, \(f\) is
  codimension-one-indeterminacy-free.
\end{corollary}
\begin{proof}
  \leanok
  A point of coheight \(1\) cannot lie in \(Z(f)\), whose points have coheight
  \(\geq 2\) by \ref{thm:indeterminacy_codimGe2}.
\end{proof}

\section{The extension theorem}
\label{sec:extension}

\begin{lemma}[Morphisms are determined by their rational-map class]
  \label{lem:ratmap_hom_ext}
  \lean{AlgebraicGeometry.Scheme.RationalMap.hom_ext_of_toRationalMap_eq}
  \leanok
  Let \(X\) be a reduced scheme, \(Y\) a separated scheme, and
  \(g_1, g_2 : X \to Y\) morphisms inducing the same rational map. Then
  \(g_1 = g_2\).
\end{lemma}
\begin{proof}
  \leanok
  Equality of the rational-map classes provides a dense open
  \(W \subseteq X\) with \(g_1|_W = g_2|_W\). The locus where \(g_1\) and
  \(g_2\) agree is the preimage of the diagonal of \(Y\) under
  \((g_1, g_2) : X \to Y \times Y\), a closed subscheme since \(Y\) is
  separated. It contains the dense open \(W\), so its underlying closed subset
  is all of \(X\); a closed subscheme of a reduced scheme with full underlying
  space is the whole scheme (its ideal sheaf consists of nilpotents). Hence
  \(g_1 = g_2\).
\end{proof}

\begin{theorem}[Extension from an empty indeterminacy locus]
  \label{thm:extension_unique_of_empty}
  \lean{AlgebraicGeometry.Scheme.RationalMap.existsUnique_hom_of_indeterminacyLocus_eq_empty}
  \uses{def:indeterminacy_locus, def:ratmap_toPartialMap, lem:ratmap_hom_ext}
  \leanok
  Let \(X\) be a reduced scheme, \(Y\) a separated scheme, and
  \(f : X \dashrightarrow Y\) a rational map with \(Z(f) = \emptyset\). Then
  there is a \emph{unique} morphism \(g : X \to Y\) whose rational map is
  \(f\).

  (The emptiness hypothesis cannot be weakened to
  codimension-one-freeness: the projection
  \(\mathbb P^2 \dashrightarrow \mathbb P^1\) away from a point has
  indeterminacy a single closed point of codimension two --- so it is
  codimension-one-free --- yet admits no regular extension.)
\end{theorem}
\begin{proof}
  \leanok
  Since \(Z(f) = \emptyset\), the maximal representative
  (\ref{def:ratmap_toPartialMap}) of \(f\) is defined on
  \(\operatorname{Dom}(f) = X\); it is a total morphism \(g\) representing
  \(f\). Uniqueness is \ref{lem:ratmap_hom_ext}.
\end{proof}

\begin{theorem}[Milne's Theorem 3.2: extension to an abelian variety]
  \label{thm:extend_to_av}
  \lean{AlgebraicGeometry.Scheme.RationalMap.extend_to_av}
  \uses{def:milne33_indeterminacy, thm:indeterminacy_codimGe2,
    thm:extension_unique_of_empty, thm:smooth_scheme_reduced}
  \leanok
  \source{abelian-varieties:page-0023}
  Let \(\bar k\) be an algebraically closed field, \(X\) a nonsingular
  \(\bar k\)-variety (as in \ref{thm:indeterminacy_codimGe2}), and \(A\) an
  abelian variety over \(\bar k\): a group scheme, proper, smooth and
  geometrically irreducible over \(\bar k\). Let \(f : X \dashrightarrow A\)
  be a rational map over \(\bar k\) which has the Milne 3.3 indeterminacy
  property (\ref{def:milne33_indeterminacy}). Then there is a \emph{unique}
  morphism \(g : X \to A\) whose rational map is \(f\).
\end{theorem}
\begin{proof}
  \leanok
  \emph{Step 0: \(A\) is a complete variety.} The scheme \(A\) is reduced by
  \ref{thm:smooth_scheme_reduced} (it is smooth over the algebraically closed
  \(\bar k\)). It is irreducible: \(\Spec \bar k\) has a single point, and a
  scheme whose structure morphism to a one-point base is geometrically
  irreducible is irreducible. Reduced and irreducible is integral; together
  with properness and geometric irreducibility, \(A\) satisfies the target
  hypotheses of \ref{thm:indeterminacy_codimGe2}. The scheme \(A\) is also
  separated (its structure morphism is separated, being proper, and
  \(\Spec \bar k\) is affine, hence separated).

  \emph{Step 1: \(Z(f) = \emptyset\).} By the Milne 3.3 property, either
  \(Z(f) = \emptyset\) and we are done, or every point of the nonempty
  \(Z(f)\) specialises from a point \(z \in Z(f)\) with
  \(\operatorname{coheight}(z) = 1\). The second alternative is impossible:
  by \ref{thm:indeterminacy_codimGe2} every point of \(Z(f)\) has coheight at
  least \(2\). Hence \(Z(f)\) has no points at all, i.e.\ it is empty.

  \emph{Step 2: extend.} \(X\) is reduced and \(A\) separated, so
  \ref{thm:extension_unique_of_empty} produces the unique morphism \(g\)
  representing \(f\).
\end{proof}
